\chapter{Conclusão}

\section{Síntese do trabalho}
\begin{figure}[H]
  \centering
  \begin{subfigure}{0.48\textwidth}
    \centering
    \includegraphics[width=\textwidth]{fotos/conclusao_inicio_modelo_video.jpg}
    \caption{Fase inicial de validação.}
  \end{subfigure}
  \hfill
  \begin{subfigure}{0.48\textwidth}
    \centering
    \includegraphics[width=\textwidth]{fotos/conclusao_final_modelo_video.jpg}
    \caption{Vista interna no ensaio final.}
  \end{subfigure}
  \caption{Comparação entre a validação inicial e o comportamento observado no modelo final.}
  \label{fig:comparacao-inicio-final-projeto}
\end{figure}

Este projeto teve como objetivo desenvolver um modelo físico e articulado de uma árvore brônquica para apoio ao treino de videobroncofibroscopia. A solução desenvolvida combinou a reconstrução anatómica a partir de dados clínicos, impressão 3D, atuação mecânica por manipuladores delta, controlo eletrónico com ESP32, comunicação com Raspberry Pi 4B, uma interface HMI e perfis de movimento inspirados em fenómenos fisiológicos observáveis durante o exame.

Ao longo do desenvolvimento, foi possível consolidar os requisitos principais do sistema e transformar uma ideia inicial ainda pouco definida num protótipo funcional, modular e capaz de reproduzir, de forma controlada, três movimentos considerados relevantes para a simulação: respiração, batimento cardíaco e tosse. A opção por manipuladores delta permitiu concentrar a atuação fora do modelo anatómico, evitando a instalação de motores diretamente na árvore brônquica e transferindo grande parte da flexibilidade do sistema para o software, através da geração de trajetórias e da cinemática inversa.

\begin{figure}[H]
  \centering
  \includegraphics[width=0.82\textwidth]{fotos/cad_sistema_final.png}
  \caption{Representação CAD do conjunto final previsto.}
  \label{fig:cad-sistema-final-conclusao}
\end{figure}

Os ensaios laboratoriais permitiram caracterizar o comportamento do manipulador e avaliar a resposta temporal dos movimentos implementados. Embora tenham sido observadas discrepâncias entre os deslocamentos pretendidos e os deslocamentos obtidos, sobretudo em amplitudes maiores e em posições mais afastadas da zona central de trabalho, os resultados mostraram que o sistema era capaz de executar movimentos coerentes com os comandos aplicados. Nos testes de frequência, tanto o batimento cardíaco como a respiração acompanharam de forma consistente os valores configurados, confirmando a viabilidade da arquitetura de controlo adotada para os perfis fisiológicos previstos.

Os testes realizados com a equipa médica permitiram ainda validar qualitativamente a utilidade do modelo enquanto ferramenta de treino. A geometria da árvore brônquica foi considerada próxima da realidade anatómica e o comportamento dinâmico introduzido pelo sistema foi visto como um contributo relevante para aproximar o simulador de um ambiente broncoscópico real. O trabalho desenvolvido cumpriu, assim, o objetivo principal de demonstrar a viabilidade de um simulador físico, acessível e evolutivo, capaz de combinar anatomia impressa em 3D com movimento controlado.

\section{Limitações}
Apesar dos resultados obtidos, o protótipo apresentou limitações importantes. A primeira esteve relacionada com o próprio modelo impresso em 3D, que não foi concebido para lavagem interna. Esta limitação reduz a capacidade de reutilização em cenários onde seja necessário aplicar lubrificantes, introduzir líquidos ou simular procedimentos que exijam limpeza posterior do interior da árvore brônquica. Além disso, a geometria atual não permite ainda a colocação simples e controlada de objetos no interior do modelo, o que limita a simulação de corpos estranhos ou de situações clínicas mais específicas.

Outra limitação relevante esteve associada à precisão dos movimentos. Os servo-motores utilizados permitiram construir uma solução compacta e de baixo custo, mas não garantiram uma correspondência perfeita entre a posição calculada e o deslocamento real da plataforma. Esta diferença resulta da resolução e repetibilidade dos atuadores, das folgas mecânicas, da fricção nas juntas impressas e das próprias aproximações usadas na cinemática inversa. Em movimentos de pequena amplitude, estas discrepâncias foram menos críticas, mas tornaram-se mais visíveis em deslocamentos maiores ou próximos dos limites mecânicos.

Foi também observada a presença de vibrações e pequenas irregularidades, sobretudo em movimentos de baixa velocidade. Estas vibrações podem introduzir inconsistências na suavidade do movimento respiratório e alterar ligeiramente a perceção externa do comportamento do modelo. Embora, durante os testes clínicos, estas vibrações se tenham revelado menos evidentes na visão interna do broncoscópio do que na observação exterior, continuam a representar uma limitação importante para futuras versões que procurem maior fidelidade e repetibilidade.

Por fim, o protótipo ainda simplifica vários fenómenos presentes numa árvore traqueobrônquica real. O modelo não reproduz deformações locais do lúmen, variações de calibre ao longo do ciclo respiratório, estenoses dinâmicas, resistência ao contacto com o broncoscópio ou alterações provocadas diretamente pela interação entre o instrumento e as paredes internas. Estes aspetos foram identificados como relevantes, mas ficaram fora do âmbito desta versão devido à complexidade mecânica, material e de controlo que implicariam.

\section{Trabalho futuro}
Como trabalho futuro, poderá ser aprofundada a evolução do modelo físico da árvore brônquica, explorando novas tecnologias e materiais de impressão que permitam obter maior flexibilidade, melhor recuperação elástica e uma resposta mecânica mais próxima dos tecidos reais. Esta evolução poderá passar pela utilização de filamentos flexíveis com diferentes durezas, resinas elásticas, silicones ou processos híbridos de fabrico, procurando combinar uma geometria anatómica detalhada com paredes mais deformáveis e resistentes ao uso repetido.

Dentro dessa evolução construtiva, poderá também ser considerada a lavagem interna do modelo. Para esse efeito, seria útil introduzir soluções que facilitem a drenagem e reduzam a acumulação de resíduos no interior dos ramos, aproximando a utilização do simulador de condições de treino mais completas e aumentando a robustez do modelo em sessões repetidas.

Outra possibilidade de desenvolvimento será a introdução de objetos ou elementos substituíveis no interior da árvore brônquica, permitindo simular corpos estranhos, obstruções ou outros cenários clínicos. Esta capacidade tornaria o modelo mais versátil, sobretudo para médicos em fase inicial de aprendizagem, ao permitir exercícios de navegação, identificação e eventual remoção de elementos colocados em posições específicas.

No sistema de atuação e controlo, poderão ser otimizados tanto o firmware do ESP32 como a interface HMI. Do lado do firmware, futuras versões poderão melhorar a suavização dos movimentos, a calibração dos atuadores, a gestão de limites mecânicos e a sincronização entre respiração, batimento cardíaco e tosse. Do lado da HMI, poderá ser desenvolvida uma interface mais intuitiva e interativa, com melhor organização dos modos de treino, maior controlo sobre os parâmetros dos movimentos, armazenamento de perfis e uma visualização mais clara do estado interno do sistema.

Por último, poderá ser vantajoso melhorar as condições de avaliação e treino disponibilizadas aos médicos. Para além da validação qualitativa já iniciada, poderão ser definidos testes e exames mais estruturados para avaliar o modelo físico, envolvendo diferentes utilizadores, níveis de experiência e cenários clínicos. Esta etapa permitiria perceber com maior rigor que componentes do simulador têm maior valor pedagógico, que limitações afetam mais a experiência de utilização e que melhorias poderiam ser priorizadas para transformar o protótipo numa ferramenta de treino mais completa e clinicamente útil.
